
\textbf{Définition 6.2.} \uline{Le foncteur défini en} 6.1 \uline{sur la catégorie des} \(K\)-\uline{algèbres} \uline{graduées}, \uline{à valeurs dans la catégorie des} \(K\)-\(\lambda\)-\uline{algèbres}, \uline{est appelé} \uline{anneau de Chern}; \uline{sa valeur pour} \(A\in\operatorname{Ob}(\mathcal G)\) \uline{est notée} \(\Ch(A)\) (\uline{ou}, \uline{s'il} \uline{y a des risques de confusion}, \(\Ch_K(A)\)).

Des formules (6.1.1), (6.1.2), (6.1.4), et (3.6.2), on déduit les relations :
\begin{equation}\tag{6.2.1}
\begin{gathered}
 (1,1+X).(1,1+Y)=(1,1+X+Y)\\
 \lambda^i(1,1+X)=0\quad\text{pour }i\ge2,
\end{gathered}
\end{equation}
relations qui servaient à définir l'anneau de Chern dans [1]. À partir de ces formules, on peut inversement retrouver les relations (6.1.1) et (6.1.4).

Indiquons la raison pour laquelle on introduit l'anneau de Chern. On se place dans la situation de l'exemple 3.9, et on désigne par \(A\) le gradué associé (par la filtration définie en 3.10) au \(K^\bullet\) d'un topos commutativement localement annelé (qu'on suppose être un \(\lambda\)-anneau pour simplifier)~; \(K\) est l'anneau d'augmentation. On verra plus loin qu'on peut définir pour tout élément d'un \(\lambda\)-anneau augmenté des classes de Chern à valeurs dans le gradué associé (6.7). À tout faisceau localement libre de type fini, on peut donc associer un élément de \(\Ch(A)\), dont la composante sur le facteur \(K\) est le rang du faisceau, et la composante sur le facteur \(1+\widehat A^+\) la somme des classes de Chern. Les formules (6.2.1) sont alors les formules usuelles qui donnent les classes de Chern et le rang du produit tensoriel de deux faisceaux inversibles et des puissances extérieures d'un faisceau inversible (c'est-à-dire que l'application qu'on vient de définir du \(K^\bullet\) dans \(\Ch(A)\) est un \(\lambda\)-homomorphisme).

